\section{Introduction}\label{sec:introduction}

Applied researchers routinely fit a single linear regression model to data pooled across naturally occurring subgroups -- age cohorts, regions, treatment arms, demographic categories -- and interpret the resulting coefficients as general descriptions of the relationship between predictors and response. This practice is common because pooled models use all available data, but it rests on the implicit assumption that the regression relationships are sufficiently similar across groups. When that assumption fails, the pooled coefficients may be poor summaries for any subgroup individually, and conclusions about sign, magnitude, and predictive accuracy may not hold outside the sample used to fit the model.

The theoretical basis for this concern is well established. Berk et al.\ \cite{berk2019} show that OLS coefficients, under misspecification, are best linear approximations to the true conditional mean function, and that these approximations depend on the marginal distribution of the regressors -- not only on the response surface. When subgroups differ in their predictor distributions, the pooled estimand itself changes, even if the true data-generating process is the same within each subgroup. Coefficient heterogeneity across groups can therefore arise from distributional differences alone, independently of any difference in underlying relationships. A well-known extreme case is Simpson's paradox, in which a pooled association reverses direction within every subgroup because the pooled estimate conflates a within-group slope with a between-group compositional effect.

A substantial literature has developed tools for detecting coefficient instability, but most of it is focused on time-ordered data. The structural break tests of Chow (1960) and the CUSUM-based procedures reviewed by Garbade \cite{garbade1977} are designed to detect discrete or gradual changes as observations accumulate over time. Watson and Engle \cite{watsenengle1985} extend this to a stationary AR(1) alternative for time-varying coefficients. These tests are powerful within their domain but assume a specific temporal structure -- a natural ordering of observations and a structured class of alternatives -- that does not apply to cross-sectional subgroup comparisons where groups are defined by observed characteristics.

A separate literature on heterogeneous treatment effects asks whether a treatment impact differs across subgroups, using methods such as interaction regression, propensity score stratification, and causal forests \cite{hte_ref}. Tsoi and Sun \cite{tsoi2026} conduct a recent Monte Carlo comparison of interaction and stratified regression for detecting subgroup differences in treatment effects, finding that the two approaches can yield conflicting conclusions depending on sample size, subgroup balance, and covariate correlation. Their results highlight the methodological risk of mixing the two approaches in reporting. That literature, however, focuses on a single binary treatment variable and frames stability as a hypothesis-testing question -- whether a subgroup difference is statistically detectable -- rather than asking how large the differences are, whether effect directions change, or when a model trained on one subgroup predicts poorly for another.

Two other lines of work approach subgroup stability from different angles. Bertsimas and Paskov \cite{bertsimas2020} propose enforcing stability by optimally selecting training and validation sets to minimize worst-case prediction error across subpopulations. Their robust optimization framework builds stability in by design and produces coefficients with lower variance and better support recovery than random splits. In panel econometrics, Liu et al.\ \cite{liu2020} develop an M-estimation approach that identifies and estimates latent group structure when subgroup membership is unknown and coefficients are heterogeneous across groups. Both papers offer solutions for settings where stability must be constructed or group membership must be recovered. Neither addresses the common diagnostic case in which groups are observed and the analyst wants to assess how stable the pooled regression conclusions already are.

This paper addresses that gap. We treat group membership as observed and ask, using a small set of interpretable metrics, how stable the conclusions of a pooled linear regression are across subpopulations in finite samples. The metrics capture four dimensions of stability: whether coefficient signs agree across groups, how much coefficient magnitudes drift relative to the pooled estimate, whether that drift is explained by differences in covariate distributions, and how much predictive accuracy degrades when a model trained on one subgroup is applied to another. We study these metrics in a controlled simulation that varies sample size, subgroup balance, predictor correlation, and the type of heterogeneity present, and we illustrate with a real-data case study using NHANES data. The goal is not to propose a new estimator but to provide practical diagnostic guidance: when should a pooled regression conclusion be treated with caution, and what does that caution look like in finite samples?

The remainder of the paper is organized as follows. Section~\ref{sec:methods} defines the stability metrics and regression models compared. Section~\ref{sec:simulation} presents the simulation study. Section~\ref{sec:realdata} presents the NHANES case study. Section~\ref{sec:discussion} discusses implications and limitations.